\documentclass[12pt,a4paper]{article}
\usepackage[margin=2.5cm]{geometry}
\usepackage[T1]{fontenc}
\usepackage[utf8]{inputenc}
\usepackage{microtype}
\usepackage{setspace}
\usepackage[colorlinks=true,urlcolor=blue]{hyperref}
\setlength{\parindent}{0pt}
\setlength{\parskip}{0.8em}

\begin{document}

3 May 2026

Editorial Office\\
\textit{OTO Open}\\
American Academy of Otolaryngology--Head and Neck Surgery Foundation

\bigskip

Dear Editor,

I submit for your consideration the manuscript titled \textbf{``Predicting
Speech-in-Noise Performance Beyond the Pure-Tone Audiogram: A Machine
Learning Benchmark on the Oldenburg Hearing Health Record''} for publication
as an Original Research article in \textit{OTO Open}.

\textbf{Relevance to OTO Open.}
Sensorineural hearing loss and its audiological management are core domains
of otolaryngology practice.  Despite widespread use of the pure-tone
audiogram, it explains only part of real-world speech understanding---leaving
clinicians with limited tools for counselling patients in the grey zone of
hearing aid candidacy.  This paper directly addresses that gap by asking:
\emph{how much speech-in-noise performance can be predicted from progressively
richer audiological measurements, and for which patients does extended
assessment add the most value?}  The answers have immediate relevance for
audiological counselling in otolaryngology clinics.

\textbf{What this paper does.}
I present the first systematic machine learning benchmark on the Oldenburg
Hearing Health Record (OHHR; Kollmeier et al., \textit{Scientific Data}
2025), an open CC~BY~4.0 dataset of 581 adults with audiograms, adaptive
loudness scaling, speech-in-noise tests, and cognitive and self-report
measures.  Five progressively richer feature blocks (1--41 features) are
evaluated across Ridge regression, Random Forest, LightGBM, and a multilayer
perceptron under 10-fold cross-validation, with SHAP-based feature
importance, severity-stratified analysis, calibrated probabilistic
classification, a minimal-predictor analysis, and partial cross-cohort
construct validation on NHANES 2017--18 ($N{=}859$).

\textbf{Key findings.}
A single binaural pure-tone average captures 71\% of speech-in-noise
variance; the full battery achieves $R^2{=}0.81$ (Random Forest,
MAE\,=\,1.23\,dB~SNR).  The clinical benefit of extended assessment is
concentrated in \emph{mild} hearing loss (26--40\,dB), where the audiogram
explains only 9\% of SIN variance and the full battery achieves 37\%---exactly
the patient group where audiogram-only counselling is most uncertain and
hearing aid candidacy most equivocal.  Age and loudness-growth slope are the
most informative features beyond the audiogram, a finding that partially
replicates in independent NHANES data (AUROC: 0.758 $\to$ 0.815 when age is
added).  Classification of poor SIN performers achieves AUROC\,=\,0.978 with
well-calibrated probability estimates within this cohort.

\textbf{Reproducibility.}
This study analyses two pre-existing, publicly available, de-identified
datasets (OHHR at Zenodo DOI 10.5281/zenodo.16919812, and NHANES via the
CDC); no institutional review board approval was required.  All analysis code
is released openly under the MIT licence.

\textbf{Suggested reviewers.}
\begin{itemize}
  \item \textbf{Prof.\ M. Samantha Lewis}, National Center for Rehabilitative
        Auditory Research, VA Portland Health Care System, USA.\\
        \textit{Expertise: clinical audiology, speech-in-noise outcomes,
        audiological rehabilitation.}

  \item \textbf{Prof.\ Michael Akeroyd}, Hearing Sciences, University of
        Nottingham, UK.
        (\href{mailto:michael.akeroyd@nottingham.ac.uk}{michael.akeroyd@nottingham.ac.uk})\\
        \textit{Expertise: hearing epidemiology, large-scale audiological
        datasets, quantitative hearing research.}

  \item \textbf{Prof.\ Tobias Dau}, Hearing Systems Group, Technical
        University of Denmark (DTU), Denmark.
        (\href{mailto:tobias.dau@dtu.dk}{tobias.dau@dtu.dk})\\
        \textit{Expertise: computational auditory modelling, speech
        intelligibility prediction, hearing science.}
\end{itemize}

\textbf{Competing interests.}
The author declares no competing financial or non-financial interests.

\textbf{AI assistance.}
Portions of the code and manuscript were drafted with the assistance of an
AI language model (Claude, Anthropic).  All scientific content, numerical
results, and interpretations were verified by the author.

I confirm that this manuscript has not been published elsewhere and is not
under consideration by any other journal.

Yours sincerely,

\bigskip
Ryan Yi Sheng Neo\\
Independent Researcher\\
\href{mailto:realryanneo@gmail.com}{realryanneo@gmail.com}\\
\url{https://github.com/ryanlearnscs50/oldenburg}

\end{document}
